\documentclass[11pt,letterpaper]{article}
\usepackage{amsmath}
\usepackage{amssymb}
\usepackage{amsthm}
\usepackage{enumitem}
\usepackage{geometry}
\usepackage{array}

\geometry{
    left=20mm,
    right=20mm,
    top=20mm,
    bottom=25mm,
}

\newcommand{\Ker}{\operatorname{Ker}}
\newcommand{\Img}{\operatorname{Im}}
\newcommand{\Tor}{\operatorname{Tor}}

\begin{document}
\pagenumbering{gobble}

\begin{center}
  Evaluasi Akhir, Semester Genap 2023/2024, Program S2, Departemen Matematika FSAD-ITS

  \vspace{0.3cm}
  \begin{tabular}{l c l}
    Matakuliah & : & Teori Modul  \\
    Tanggal    & : & 19 Juni 2024 \\
    Waktu      & : & 100 menit    \\
    Sifat      & : & Tutup buku   \\
    Dosen      & : & Subiono
  \end{tabular}

  \vspace{0.3cm}
  \rule{0.45\textwidth}{0.4pt}
\end{center}

\vspace{0.5cm}

\begin{enumerate}
  \item Bila $M$ adalah suatu modul-$R$ dan bila $A,B,C$ adalah submodul $M$ dengan $C \subseteq A$, maka tunjukkan bahwa
        $
          \boxed{A \cap (B+C) = (A \cap B) + C}.
        $

  \item Misalkan $M$ adalah suatu modul atas suatu ring $R$. Bila $\underline{n}=\{1,2,3,\ldots,n\}$ dan
        $
          M = \bigoplus_{i \in \underline{n}} M_i
        $
        dengan $M_i$ adalah submodul bebas dari $M$ untuk semua $i \in \underline{n}$. Maka tunjukkan bahwa $M$ adalah modul bebas.

  \item Misalkan $M$ adalah suatu modul atas suatu Daerah Ideal Utama $R$ yang dibangun secara berhingga. Bila $M \neq M_t$ dengan $M_t$ adalah modul torsi, maka tunjukkan bahwa ada submodul bebas $N$ dari $M$ yang memenuhi
        $
          M = M_t \oplus N.
        $

  \item Bila
        \[
          M_1 \xrightarrow{f_1} M_2 \xrightarrow{f_2} M_3 \xrightarrow{f_3} M_4
        \]
        adalah barisan eksak dari modul atas ring $R$. Maka tunjukkan bahwa dua barisan berikut
        \[
          M_1 \xrightarrow{f_1} M_2 \xrightarrow{f_2} N \to 0
          \qquad \text{dan} \qquad
          0 \to N \xrightarrow{i} M_3 \xrightarrow{f_3} M_4
          \tag{$A$}
        \]
        juga merupakan barisan eksak, dimana $N = \Img(f_2) = \Ker(f_3)$ dan $i$ adalah pemetaan inklusi.

        Sebaliknya bila dua barisan
        \[
          M_1 \xrightarrow{f_1} M_2 \xrightarrow{f_2} N \to 0
          \qquad \text{dan} \qquad
          0 \to N \xrightarrow{i} M_3 \xrightarrow{f_3} M_4
          \tag{$A$}
        \]
        adalah barisan eksak dengan $N \leq M_3$ dan $i$ adalah pemetaan inklusi, maka tunjukkan bahwa barisan
        \[
          M_1 \xrightarrow{f_1} M_2 \xrightarrow{f_2} M_3 \xrightarrow{f_3} M_4
          \tag{$B$}
        \]
        adalah eksak.
\end{enumerate}

\newpage
\begin{center}
  \textbf{SOLUSI}
\end{center}

\begin{enumerate}
  \item Akan ditunjukkan dua inklusi.

        Pertama, ambil $x \in A \cap (B+C)$. Maka $x \in A$ dan $x \in B+C$, sehingga terdapat $b \in B$ dan $c \in C$ dengan
        \[
          x=b+c.
        \]
        Karena $C \subseteq A$, maka $c \in A$. Selain itu $x \in A$, sehingga
        \[
          b=x-c \in A.
        \]
        Karena juga $b \in B$, diperoleh $b \in A \cap B$. Jadi
        \[
          x=b+c \in (A \cap B)+C.
        \]
        Maka
        \[
          A \cap (B+C) \subseteq (A \cap B)+C.
        \]

        Sebaliknya, ambil $x \in (A \cap B)+C$. Maka terdapat $y \in A \cap B$ dan $c \in C$ sehingga
        \[
          x=y+c.
        \]
        Karena $y \in A$ dan $c \in C \subseteq A$, maka $x \in A$. Karena $y \in B$ dan $c \in C$, maka $x \in B+C$. Jadi
        \[
          x \in A \cap (B+C).
        \]
        Maka
        \[
          (A \cap B)+C \subseteq A \cap (B+C).
        \]
        Dari dua inklusi tersebut,
        \[
          \boxed{A \cap (B+C) = (A \cap B)+C}.
        \]

  \item Untuk setiap $i \in \underline{n}$, karena $M_i$ adalah modul bebas, ambil basis $S_i$ dari $M_i$. Akan ditunjukkan bahwa
        \[
          S=\bigcup_{i=1}^{n} S_i
        \]
        adalah basis bagi $M$.

        Karena
        \[
          M=\bigoplus_{i=1}^{n} M_i,
        \]
        setiap $m \in M$ dapat ditulis secara tunggal sebagai
        \[
          m=m_1+m_2+\cdots+m_n,
        \]
        dengan $m_i \in M_i$. Karena $S_i$ basis dari $M_i$, setiap $m_i$ merupakan kombinasi linear berhingga dari elemen-elemen $S_i$. Maka $m$ merupakan kombinasi linear berhingga dari elemen-elemen $S$. Jadi $S$ membangun $M$.

        Selanjutnya, misalkan terdapat relasi linear berhingga
        \[
          r_1s_1+r_2s_2+\cdots+r_ks_k=0,
        \]
        dengan $r_j \in R$ dan $s_j \in S$. Kelompokkan suku-suku tersebut berdasarkan submodul $M_i$. Maka relasi itu dapat ditulis sebagai
        \[
          u_1+u_2+\cdots+u_n=0,
        \]
        dengan $u_i \in M_i$ dan masing-masing $u_i$ merupakan kombinasi linear dari elemen-elemen $S_i$. Karena jumlahnya adalah jumlah langsung, representasi nol bersifat tunggal, sehingga
        \[
          u_i=0
          \quad \text{untuk setiap } i.
        \]
        Karena $S_i$ bebas linear di $M_i$, semua koefisien pada $u_i$ bernilai nol. Jadi semua $r_j=0$. Maka $S$ bebas linear.

        Dengan demikian $S$ membangun $M$ dan bebas linear, sehingga $S$ adalah basis dari $M$. Jadi
        \[
          \boxed{M \text{ adalah modul bebas}.}
        \]

  \item Misalkan
        \[
          M_t=\{m \in M \mid \text{terdapat }0 \neq r \in R \text{ sehingga } rm=0\}
        \]
        adalah submodul torsi dari $M$. Karena $M$ dibangun secara berhingga atas Daerah Ideal Utama $R$, maka $M/M_t$ juga dibangun secara berhingga. Selain itu $M/M_t$ adalah modul bebas torsi. Memang, jika
        \[
          0 \neq r \in R
          \quad \text{dan} \quad
          r(m+M_t)=M_t,
        \]
        maka $rm \in M_t$. Jadi terdapat $0 \neq s \in R$ sehingga
        \[
          s(rm)=0.
        \]
        Karena $R$ domain integral, $sr \neq 0$, sehingga $m \in M_t$. Maka $m+M_t=M_t$, sehingga $M/M_t$ bebas torsi.

        Karena modul yang dibangun secara berhingga dan bebas torsi atas Daerah Ideal Utama adalah modul bebas, maka $M/M_t$ bebas. Karena $M \neq M_t$, basis $M/M_t$ tidak kosong. Misalkan
        \[
          \{x_1+M_t,x_2+M_t,\ldots,x_k+M_t\}
        \]
        adalah basis dari $M/M_t$, dengan $x_i \in M$. Definisikan
        \[
          N=Rx_1+Rx_2+\cdots+Rx_k.
        \]

        Akan ditunjukkan bahwa $N$ bebas. Jika
        \[
          a_1x_1+a_2x_2+\cdots+a_kx_k=0,
        \]
        maka dalam $M/M_t$ berlaku
        \[
          a_1(x_1+M_t)+a_2(x_2+M_t)+\cdots+a_k(x_k+M_t)=M_t.
        \]
        Karena $\{x_i+M_t\}_{i=1}^{k}$ basis, diperoleh $a_1=a_2=\cdots=a_k=0$. Jadi $N$ bebas dengan basis $\{x_1,x_2,\ldots,x_k\}$.

        Selanjutnya, untuk setiap $m \in M$, karena $\{x_i+M_t\}_{i=1}^{k}$ membangun $M/M_t$, terdapat $a_1,\ldots,a_k \in R$ sehingga
        \[
          m+M_t=a_1(x_1+M_t)+\cdots+a_k(x_k+M_t).
        \]
        Maka
        \[
          m-(a_1x_1+\cdots+a_kx_k) \in M_t.
        \]
        Jadi $m \in M_t+N$, sehingga
        \[
          M=M_t+N.
        \]

        Terakhir, jika $y \in M_t \cap N$, maka
        \[
          y=a_1x_1+\cdots+a_kx_k
        \]
        untuk suatu $a_1,\ldots,a_k \in R$. Karena $y \in M_t$, dalam $M/M_t$ diperoleh
        \[
          a_1(x_1+M_t)+\cdots+a_k(x_k+M_t)=M_t.
        \]
        Karena basis bebas linear, semua $a_i=0$, sehingga $y=0$. Jadi
        \[
          M_t \cap N=\{0\}.
        \]

        Dengan demikian
        \[
          \boxed{M=M_t \oplus N}
        \]
        dengan $N$ submodul bebas dari $M$.

  \item Misalkan
        \[
          M_1 \xrightarrow{f_1} M_2 \xrightarrow{f_2} M_3 \xrightarrow{f_3} M_4
        \]
        eksak. Artinya
        \[
          \Img(f_1)=\Ker(f_2)
          \quad \text{dan} \quad
          \Img(f_2)=\Ker(f_3).
        \]
        Ambil
        \[
          N=\Img(f_2)=\Ker(f_3)
        \]
        dan $i:N \to M_3$ pemetaan inklusi.

        Untuk barisan pertama
        \[
          M_1 \xrightarrow{f_1} M_2 \xrightarrow{f_2} N \to 0,
        \]
        pemetaan $f_2$ dipandang sebagai pemetaan dari $M_2$ ke $N=\Img(f_2)$. Maka $f_2:M_2 \to N$ surjektif. Jadi
        \[
          \Img(f_2)=N=\Ker(N \to 0).
        \]
        Selain itu, karena barisan semula eksak,
        \[
          \Img(f_1)=\Ker(f_2).
        \]
        Maka
        \[
          M_1 \xrightarrow{f_1} M_2 \xrightarrow{f_2} N \to 0
        \]
        adalah eksak.

        Untuk barisan kedua
        \[
          0 \to N \xrightarrow{i} M_3 \xrightarrow{f_3} M_4,
        \]
        karena $i$ adalah inklusi, maka $i$ injektif, sehingga
        \[
          \Img(0 \to N)=0=\Ker(i).
        \]
        Selain itu
        \[
          \Img(i)=N=\Ker(f_3).
        \]
        Jadi
        \[
          0 \to N \xrightarrow{i} M_3 \xrightarrow{f_3} M_4
        \]
        adalah eksak. Dengan demikian dua barisan pada (A) eksak.

        Sebaliknya, misalkan dua barisan
        \[
          M_1 \xrightarrow{f_1} M_2 \xrightarrow{g} N \to 0
          \qquad \text{dan} \qquad
          0 \to N \xrightarrow{i} M_3 \xrightarrow{f_3} M_4
        \]
        eksak, dengan $N \leq M_3$ dan $i$ pemetaan inklusi. Karena barisan pertama eksak,
        \[
          \Img(f_1)=\Ker(g)
          \quad \text{dan} \quad
          g \text{ surjektif ke } N.
        \]
        Karena barisan kedua eksak,
        \[
          i \text{ injektif}
          \quad \text{dan} \quad
          \Img(i)=N=\Ker(f_3).
        \]
        Definisikan pemetaan pada barisan (B) dari $M_2$ ke $M_3$ sebagai
        \[
          f_2=i \circ g.
        \]
        Karena $i$ inklusi, $\Ker(i)=0$, sehingga
        \[
          \Ker(f_2)
          =\Ker(i \circ g)
          =\Ker(g)
          =\Img(f_1).
        \]
        Jadi barisan (B) eksak di $M_2$.

        Selanjutnya,
        \[
          \Img(f_2)
          =\Img(i \circ g)
          =i(\Img(g)).
        \]
        Karena $g$ surjektif ke $N$, maka $\Img(g)=N$. Jadi
        \[
          \Img(f_2)=i(N)=N.
        \]
        Dari eksaknya barisan kedua, $N=\Ker(f_3)$. Maka
        \[
          \Img(f_2)=\Ker(f_3).
        \]
        Jadi barisan (B) eksak di $M_3$.

        Dengan demikian
        \[
          \boxed{
            M_1 \xrightarrow{f_1} M_2 \xrightarrow{f_2} M_3 \xrightarrow{f_3} M_4
            \text{ adalah barisan eksak}.
          }
        \]
\end{enumerate}

\end{document}
